\begin{table}[!htbp]
\centering
\caption[LSTM vs predictori naivi]{Validarea LSTM față de predictori naivi, eroare o-zi-înainte pe fereastra de test (ultimele $20\%$ din seria România, $288$ zile). Toate metodele folosesc doar valori observate trecute. \emph{Skill} este reducerea relativă a MAE față de persistența (valori negative = mai slab decât persistența). Pe această serie, $85\%$ din zilele de test sunt identice cu ziua precedentă (raportare săptămânală), ceea ce face persistența extrem de greu de bătut.}
\label{tab:lstm_baselines}
\small
\begin{tabularx}{\columnwidth}{X r r r}
\toprule
\textbf{Metodă} & \textbf{MAE} & \textbf{RMSE} & \textbf{Skill} \\
\midrule
LSTM (o-zi-înainte) & 755 & 1114 & -200\% \\
Persistență & 251 & 963 & --- \\
Drift (ultima diferență) & 501 & 1362 & -99\% \\
Medie mobilă (7 zile) & 1002 & 1628 & -298\% \\
\bottomrule
\end{tabularx}
\end{table}
